% Notação vetorial; não depende de PlantUML, Java ou serviços externos.
\usetikzlibrary{shapes.geometric,shapes.multipart,calc}
\tikzset{
  uml/caixa/.style={draw=srslinha,fill=white,align=center,font=\small,inner sep=5pt},
  uml/caso/.style={uml/caixa,ellipse,text width=38mm,minimum height=17mm},
  uml/classe/.style={uml/caixa,rectangle split,rectangle split parts=3,
    rectangle split part align={center,left,left},text width=48mm},
  uml/componente/.style={uml/caixa,rectangle,text width=35mm,minimum height=18mm},
  uml/acao/.style={uml/caixa,rounded corners=4pt,text width=43mm,minimum height=10mm},
  uml/decisao/.style={uml/caixa,diamond,aspect=2.5,text width=25mm,inner sep=1pt},
  uml/estado/.style={uml/caixa,rounded corners=6pt,fill=srscabecalho,
    minimum width=30mm,minimum height=13mm},
  uml/fluxo/.style={-{Stealth[length=2mm]},semithick},
  uml/chamada/.style={-{Triangle[length=2mm,width=1.7mm]},semithick},
  uml/retorno/.style={dashed,-{Straight Barb[length=2mm]},semithick},
  uml/dependencia/.style={dashed,-{Straight Barb[length=2mm]},semithick},
  uml/rotulo/.style={font=\footnotesize,align=center,inner sep=2pt},
  uml/nota/.style={draw=srslinha,dashed,align=left,font=\footnotesize,
    fill=white,text width=55mm,inner sep=5pt}
}

\newcommand{\atoruml}[3]{%
  \begin{scope}[shift={#2}]
    \node[minimum width=8mm,minimum height=15mm,inner sep=0pt] (#1) at (0,0) {};
    \draw (0,0.48) circle (0.16);
    \draw (0,0.32) -- (0,-0.22);
    \draw (-0.32,0.12) -- (0.32,0.12);
    \draw (-0.27,-0.55) -- (0,-0.22) -- (0.27,-0.55);
    \node[uml/rotulo,text width=25mm] at (0,-0.92) {#3};
  \end{scope}}
\newcommand{\finaluml}[2]{%
  \node[draw,circle,minimum size=5mm,inner sep=0pt] (#1) at #2 {};
  \fill (#1.center) circle (1.4mm);}
\newcommand{\pacoteuml}[3]{%
  \node[uml/caixa,minimum width=45mm,minimum height=15mm] (#1) at #2 {#3};
  \draw[draw=srslinha] (#1.north west) -- ++(0,0.4) -- ++(1.5,0) -- ++(0,-0.4);}
\newcommand{\nouml}[4]{%
  \node[uml/caixa,text width=40mm,minimum height=36mm] (#1) at #2 {};
  \node[font=\small,align=center,text width=43mm,anchor=north]
    at ($(#1.north)+(0,-0.16)$) {#3};
  \node[uml/caixa,text width=36mm,inner sep=3pt,anchor=south]
    at ($(#1.south)+(0,0.16)$) {#4};
  \draw[draw=srslinha] (#1.north west) -- ++(0.3,0.3)
    -- ($(#1.north east)+(0.3,0.3)$) -- (#1.north east);
  \draw[draw=srslinha] (#1.north east) -- ++(0.3,0.3)
    -- ($(#1.south east)+(0.3,0.3)$) -- (#1.south east);}
\newcommand{\diagramasrs}[3]{%
  \begin{figurasrs}
  \caption{#1}\label{#2}
  \centering
  #3\par
  \fonte{Elaboração didática deste guia.}
  \end{figurasrs}}
